\begin{frame}{Argumento de troca}
    \metroset{block=fill}

    \begin{block}{Método do argumento de troca}
        O argumento de troca é muito utilizado para mostrar que certas escolhas não precisam ser consideradas. \pause
        \footnotesize
        \begin{itemize}
            \item {\bf Passo 1}. Consideramos uma possível resposta. \pause
            \item {\bf Passo 2}. Mostramos que, a partir de uma certa troca de escolhas, a nova escolha continua sendo uma resposta. \pause
            \item {\bf Passo 3}. Analisamos quais escolhas nunca valem a pena ser consideradas.
        \end{itemize}
    \end{block} \pause
    \footnotesize
    Analisemos novamente a prova do slide anterior: \pause
    \begin{itemize}
        \item Passo 1. Consideramos um subconjunto $S$ com a maior quantidade possível de algoritmos. \pause
        \item Passo 2. Se trocarmos um algoritmo $i \in S$ por um algoritmo $j \in S$ tal que $a_i > a_j$, o novo conjunto continua sendo uma resposta válida. \pause
        \item Passo 3. Concluimos que não vale a pena considerar os subconjuntos que não consistem nos menores elementos, já que sempre podemos trocar um algoritmo em $S$ por outro com tempo menor fora de $S$.
    \end{itemize}
    % Note que só existe uma única resposta que consiste apenas nos menores elementos.

\end{frame}

